%================================
% preamble.tex — 词汇讲义公用导言区
% 基于 note-setup-borderless.tex (fenglielie@qq.com 2025-09-12) 改名并扩充
%================================

\usepackage{amsmath,amsthm,amsfonts,amssymb}
\usepackage{mathtools}
\usepackage{mathrsfs}
\usepackage{bm}
\usepackage{extarrows}
\usepackage[a4paper, margin=1in]{geometry}
\usepackage{float}
\usepackage{indentfirst}
\usepackage{anyfontsize}
\usepackage{booktabs,multirow,multicol}
\usepackage[shortlabels,inline]{enumitem}
\usepackage{appendix}
\usepackage[UTF8]{ctex}
\usepackage{fontspec}
\usepackage{booktabs}

\geometry{a4paper,left=1.5cm,right=1.5cm,top=1.5cm,bottom=1.5cm}
\setmainfont{Times New Roman} % 设置正文主字体
\setsansfont{Arial} % 设置无衬线字体
\setmonofont{Courier New} % 设置等宽字体
\setCJKmainfont{SimSun}[AutoFakeBold=2.0] % 设置中文主字体（宋体）
\setCJKsansfont{SimHei} % 设置中文无衬线字体（黑体）
\setCJKmonofont{FangSong} % 设置中文等宽字体（仿宋）
\pagestyle{plain}
\setcounter{tocdepth}{1}

\usepackage[dvipsnames,svgnames]{xcolor}
\usepackage{graphicx}
\graphicspath{
    {./figure/}{./figures/}{./image/}{./images/}{./graphic/}{./graphics/}{./picture/}{./pictures/}
}
\usepackage{subcaption}

\usepackage[ruled,linesnumbered,noline]{algorithm2e}
\usepackage{listings}
\lstdefinestyle{simpleStyle}{
    basicstyle=\ttfamily\small,
    breaklines=true,
    keywordstyle=\color{blue},
    identifierstyle=\color{black},
    stringstyle=\color{violet},
    commentstyle=\color[RGB]{34,139,34},
    showstringspaces=false,
    numbers=left,
    numbersep=2em,
    numberstyle=\footnotesize,
    frame=single,
    framesep=1em,
}
\lstset{style=simpleStyle}

\usepackage{hyperref}
\hypersetup{
    colorlinks=true,linkcolor=,urlcolor=cyan
}

\renewcommand*{\proofname}{\normalfont\bfseries Proof}

\usepackage{thmtools}

%% define environments

\declaretheorem[style=plain, name=Theorem, numbered=yes, numberwithin=section]{theorem}
\declaretheorem[style=plain, name=Theorem, numbered=no]{theorem*}

\declaretheorem[style=plain, name=Proposition, numbered=yes, sibling=theorem]{proposition}
\declaretheorem[style=plain, name=Proposition, numbered=no]{proposition*}

\declaretheorem[style=plain, name=Law, numbered=yes, sibling=theorem]{law}
\declaretheorem[style=plain, name=Law, numbered=no]{law*}

\declaretheorem[style=plain, name=Formula, numbered=yes, sibling=theorem]{formula}
\declaretheorem[style=plain, name=Formula, numbered=no]{formula*}

% \declaretheorem[style=plain, name=Corollary, numbered=yes, sibling=theorem]{corollary}
% \declaretheorem[style=plain, name=Corollary, numbered=no]{corollary*}

\declaretheorem[style=plain, name=Lemma, numbered=yes, sibling=theorem]{lemma}
\declaretheorem[style=plain, name=Lemma, numbered=no]{lemma*}

\declaretheorem[style=plain, name=Claim, numbered=yes, sibling=theorem]{claim}
\declaretheorem[style=plain, name=Claim, numbered=no]{claim*}

\declaretheorem[style=definition, name=Definition, numbered=yes, numberwithin=section]{definition}
\declaretheorem[style=definition, name=Definition, numbered=no]{definition*}

\declaretheorem[style=definition, name=Example, numbered=yes, numberwithin=section]{example}
\declaretheorem[style=definition, name=Example, numbered=no]{example*}

\declaretheorem[style=definition, name=Problem, numbered=yes, numberwithin=section]{problem}
\declaretheorem[style=remark, name=Remark, numbered=yes, numberwithin=section]{remark}
\declaretheorem[style=remark, name=Remark, numbered=no]{remark*}

\declaretheoremstyle[headfont=\color{orange!80}\bfseries, bodyfont=\normalfont, spaceabove=3pt, spacebelow=3pt]{notestyle}

\declaretheorem[style=notestyle, name=Note, numbered=yes, numberwithin=section]{note}
\declaretheorem[style=notestyle, name=Note, numbered=no]{note*}

\declaretheoremstyle[headfont=\bfseries, bodyfont=\normalfont, spaceabove=3pt, spacebelow=3pt, qed=\ensuremath{\square}]{solutionstyle}

\declaretheorem[style=solutionstyle, name=Solution, numbered=yes, numberwithin=section]{solution}
\declaretheorem[style=solutionstyle, name=Solution, numbered=no]{solution*}

% thmtools numberwithin=section does not always propagate the section prefix to
% hyperref destinations. Include the section anchor explicitly to avoid duplicate
% PDF destination names such as example.1 in different sections.
\makeatletter
\@ifpackageloaded{hyperref}{%
    \newcommand{\fixTheoremHref}[1]{%
        \@namedef{theH#1}{\theHsection.\arabic{#1}}%
    }%
    \fixTheoremHref{theorem}%
    \fixTheoremHref{definition}%
    \fixTheoremHref{example}%
    \fixTheoremHref{problem}%
    \fixTheoremHref{remark}%
    \fixTheoremHref{note}%
    \fixTheoremHref{solution}%
    \@namedef{theHproposition}{\theHtheorem}%
    \@namedef{theHcorollary}{\theHtheorem}%
    \@namedef{theHlemma}{\theHtheorem}%
    \@namedef{theHclaim}{\theHtheorem}%
}{}
\makeatother

\usepackage[most]{tcolorbox}

\newcommand{\newtcbenvironment}[2]{
    \tcolorboxenvironment{#1}{#2, enhanced, breakable, sharp corners, boxrule=0pt, colframe=white}
    \tcolorboxenvironment{#1*}{#2, enhanced, breakable, rounded corners, boxrule=0pt, colframe=white}
}

%% define styles

\newtcbenvironment{theorem}{colframe=RoyalPurple, colback=RoyalPurple!8}
\newtcbenvironment{proposition}{colframe=RoyalPurple, colback=RoyalPurple!8}
\newtcbenvironment{law}{colframe=RoyalPurple, colback=RoyalPurple!8}
\newtcbenvironment{formula}{colframe=RoyalPurple, colback=RoyalPurple!8}
\newtcbenvironment{corollary}{colframe=NavyBlue, colback=SkyBlue!8}
\newtcbenvironment{lemma}{colframe=NavyBlue, colback=SkyBlue!8}
\newtcbenvironment{claim}{colframe=NavyBlue, colback=SkyBlue!8}

\newtcbenvironment{definition}{colframe=ForestGreen, colback=ForestGreen!5}
\newtcbenvironment{example}{colframe=RawSienna, colback=RawSienna!5}
\newtcbenvironment{problem}{colframe=WildStrawberry!30, colback=WildStrawberry!5}

%% cbox
\newtcolorbox{cbox}[1][]{%
    enhanced,
    breakable,
    sharp corners,
    leftrule=2pt, rightrule=0pt, toprule=0pt, bottomrule=0pt,
    colframe=SkyBlue,
    colback=SkyBlue!8,
    #1
}

%% cover
\usepackage{titling}
\newcommand{\extrainfo}{}
\renewcommand{\extrainfo}[1]{\renewcommand{\extrainfocontent}{#1}}
\newcommand{\extrainfocontent}{}
\newcommand{\makecover}[1]{%
    \begin{titlepage}
        \newgeometry{margin=0in}
        \parindent=0pt
        \includegraphics[width=\linewidth]{#1} % size = 1280*1024
        \vfill
        \begin{center}
            \parbox{0.618\textwidth}{%
                \raggedleft{\bfseries \fontsize{40}{3}\selectfont \thetitle} \\[0.6pt]
                \rule{0.618\textwidth}{4pt} \\
            }
        \end{center}
        \vfill
        \begin{center}
            \parbox{0.618\textwidth}{%
                \raggedleft\Large
                \begin{tabular}{r}
                    \theauthor \\
                    \thedate   \\
                \end{tabular}%
            }
        \end{center}
        \vfill
        \begin{center}
            \parbox[t]{0.7\textwidth}{\centering \itshape \extrainfocontent}
        \end{center}
        \vfill
    \end{titlepage}
    \restoregeometry
    \thispagestyle{empty}
}

\renewcommand{\vec}{\math}

%================================
% 词汇讲义专用样式
%================================
\usepackage{pifont}

% 词条标题: \vw{单词}{音标}{词性速览}
\newcommand{\vw}[3]{%
  \par\bigskip\noindent
  {\color{SkyBlue}\bfseries\Large #1}\
  {\small\itshape [{#3}]}\
  {\small\color{gray} {#2}}%
  \par\smallskip
}
% 短语词条标题: \vp{短语}
\newcommand{\vp}[1]{%
  \par\bigskip\noindent
  {\color{NavyBlue}\bfseries\Large #1}%
  \par\smallskip
}
% 义项: \sense{词性}{中文释义}
\newcommand{\sense}[2]{\par\medskip\noindent{\bfseries\color{RoyalPurple}[{#1}]}\ {\bfseries {#2}}\par\smallskip}
% 搭配/用法行: \col{搭配或句型}
\definecolor{myBurgundy}{RGB}{128,0,32}
\newcommand{\col}[1]{\par\noindent\hspace{1em}{\color{myBurgundy}$\blacktriangleright$\ \bfseries {#1}}\par\nopagebreak\smallskip}
% 例句(广义,含短语级示例): \ex{英文}
\newcommand{\ex}[1]{\par\nopagebreak\noindent\hspace{2.2em}{\itshape {#1}}\par\smallskip}
% 语法侧栏
\newtcolorbox{gram}[1]{enhanced, breakable, sharp corners,
  leftrule=2.5pt, rightrule=0pt, toprule=0pt, bottomrule=0pt,
  colframe=RawSienna, colback=RawSienna!5,
  title={\bfseries Language Focus: {#1}}, fonttitle=\small,
  coltitle=white, colbacktitle=orange!80!black,
  attach boxed title to top left={yshift=-2mm, xshift=3mm},
  boxed title style={sharp corners, boxrule=0pt},
  top=4mm}
% 辨析侧栏
\newtcolorbox{discbox}[1]{enhanced, breakable, sharp corners,
  leftrule=2.5pt, rightrule=0pt, toprule=0pt, bottomrule=0pt,
  colframe=TealBlue, colback=TealBlue!6,
  title={\bfseries Lexical discrimination: {#1}}, fonttitle=\small,
  coltitle=white, colbacktitle=TealBlue,
  attach boxed title to top left={yshift=-2mm, xshift=3mm},
  boxed title style={sharp corners, boxrule=0pt},
  top=4mm}
% 词源侧栏
\newtcolorbox{etymbox}{enhanced, breakable, sharp corners,
  leftrule=2.5pt, rightrule=0pt, toprule=0pt, bottomrule=0pt,
  colframe=ForestGreen, colback=ForestGreen!5,
  title={\bfseries Word Origin}, fonttitle=\small,
  coltitle=white, colbacktitle=OliveGreen,
  attach boxed title to top left={yshift=-2mm, xshift=3mm},
  boxed title style={sharp corners, boxrule=0pt},
  top=4mm}
